% =========================================================
% Preàmbul comú per als fragments d'exercicis (enunciat.tex / solucio.tex).
% S'inclou DESPRÉS de \documentclass i babel. El fan servir comprova.py
% (previsualització) i, més endavant, l'app per compondre exàmens.
% =========================================================
\usepackage{amsmath,amssymb}
\usepackage{graphicx}
\usepackage{xcolor}
\usepackage{icomma}        % $6,87$ amb coma decimal i espaiat correcte
\usepackage{enumitem}
\usepackage{array}
\usepackage{booktabs}
\usepackage{newunicodechar}
% "l·l" ben compost amb el caràcter UTF-8 (i "·" dins de fórmules = \cdot)
\newunicodechar{·}{\ifmmode\cdot\else\kern-0.06em\raisebox{0.42ex}{.}\kern-0.06em\fi}

% Carpeta de l'exercici que s'està incloent (la redefineix qui l'inclou).
\providecommand{\exdir}{.}

% ---------- Figures ----------
% \figura[amplada relativa]{fitxer.png}
\newcommand{\figura}[2][0.5]{%
  \par\smallskip{\centering\includegraphics[width=#1\linewidth]{\exdir/#2}\par}\smallskip}

% ---------- Enunciat ----------
% Bloc de dades i de notes al final de l'enunciat.
\newcommand{\blocetiquetat}[2]{%
  \par\medskip\noindent
  \makebox[1.6cm][l]{\textsc{#1}}%
  \parbox[t]{\dimexpr\linewidth-1.6cm\relax}{\raggedright #2}\par\smallskip}
\newcommand{\dades}[1]{\blocetiquetat{Dades:}{#1}}
\newcommand{\nota}[1]{\blocetiquetat{Nota:}{#1}}

% Apartats: \begin{apartat}{1,25} ... \end{apartat}
% L'argument són els punts (text, amb coma decimal). L'etiqueta la posa qui
% compon el document (\etiquetaapartat), per poder renumerar.
\newcounter{apartat}
\providecommand{\etiquetaapartat}{\alph{apartat})}
% \mostrapuntsfalse amaga els punts dels apartats (fulls d'exercicis).
\newif\ifmostrapunts \mostrapuntstrue
% "1 punt" / "1,25 punts"
\newcommand{\textpunts}[1]{#1~\def\puntsu{1}\edef\puntsx{#1}\ifx\puntsx\puntsu punt\else punts\fi}
% Posa #1 alineat a la dreta al final del paràgraf (a la línia següent si no hi cap).
\newcommand{\aladreta}[1]{{\unskip\nobreak\hfil\penalty50\hskip1em\hbox{}\nobreak\hfil
  \mbox{#1}\parfillskip=0pt \finalhyphendemerits=0 \par}}
% És una llista d'un sol element: els paràgrafs interns conserven el sagnat.
\newenvironment{apartat}[1]{%
  \def\puntsapartat{#1}%
  \refstepcounter{apartat}%
  \begin{list}{}{\setlength{\leftmargin}{2em}\setlength{\labelwidth}{1.6em}%
    \setlength{\labelsep}{0.4em}\setlength{\topsep}{4pt}%
    \setlength{\parsep}{3pt}\setlength{\itemsep}{0pt}}%
  \item[\textbf{\etiquetaapartat}]\ignorespaces
}{%
  \ifmostrapunts\ifx\puntsapartat\empty\else\aladreta{\textit{[\textpunts{\puntsapartat}]}}\fi\fi
  \end{list}}

% ---------- Solució ----------
% \begin{solapartat} ... \end{solapartat}  (un per apartat, mateix ordre)
\newcounter{solapartat}
\providecommand{\etiquetasolapartat}{\alph{solapartat})}
\newenvironment{solapartat}{%
  \refstepcounter{solapartat}%
  \par\medskip\noindent\textbf{\etiquetasolapartat}\par\nopagebreak\smallskip
}{\par}
% Puntuació parcial de la pauta: \punts{0,75}
\newcommand{\punts}[1]{\textcolor{red!70!black}{\textbf{#1\,p}}\ }
% Paraula destacada de la pauta ("Alternativament", "Nota"...)
\newcommand{\destaca}[1]{\textcolor{blue!70!black}{\textbf{#1}}}
